%% sections/study_objects.tex
\section{An in-depth study of the various features of the Drum Tower}
\label{sec:study_objects}

\subsubsection{\texorpdfstring{\textbf{\hl{3 Study
Objects}}}{3 Study Objects}}\label{study-objects}

\hl{The four selected cases---Chaoli, Congjiang, Zeli, and
Zengchong---provide a representative empirical basis for method
comparison. The external forms and structural characteristics of the
four drum towers are shown in Fig.5. This selection strategy emphasizes
heterogeneity in reconstruction difficulty, covering variations in roof
hierarchy, timber density, shadowed spaces, viewpoint continuity, and
village surroundings.}

\includegraphics[width=6.6875in,height=4.41181in]{figures/media/image7.png}

\textbf{\hl{Fig. 5. Four Dong drum tower cases.}}

\hl{As shown in Table 2, Congjiang Drum Tower has a relatively clear
external geometry, but still presents challenges in handling repetitive
timber textures and preserving local details. Zeli Drum Tower is
characterized by a distinctly layered roof structure, making it suitable
for evaluating the ability of different methods to reconstruct repeated
eave patterns and subtle depth relationships between roof tiers. Chaoli
Drum Tower represents a more complex acquisition scenario, with poor
viewpoint continuity and unstable test-view construction conditions,
which directly affect view synthesis quality and cross-method
comparability. Zengchong Drum Tower contains dense multi-layered eave
geometry and rich timber-component details, posing substantial
challenges for eva.}

\begin{longtable}[]{@{}
  >{\centering\arraybackslash}p{(\linewidth - 8\tabcolsep) * \real{0.0654}}
  >{\centering\arraybackslash}p{(\linewidth - 8\tabcolsep) * \real{0.1289}}
  >{\centering\arraybackslash}p{(\linewidth - 8\tabcolsep) * \real{0.2499}}
  >{\centering\arraybackslash}p{(\linewidth - 8\tabcolsep) * \real{0.3237}}
  >{\centering\arraybackslash}p{(\linewidth - 8\tabcolsep) * \real{0.2321}}@{}}
\caption{\textbf{\hl{Table 2. Overview of the four Dong drum tower cases
and their main reconstruction challenges}}}\tabularnewline
\toprule\noalign{}
\begin{minipage}[b]{\linewidth}\centering
\textbf{\hl{Case ID}}
\end{minipage} & \begin{minipage}[b]{\linewidth}\centering
\textbf{\hl{Scene}}
\end{minipage} & \begin{minipage}[b]{\linewidth}\centering
\textbf{\hl{Architectural characteristics}}
\end{minipage} & \begin{minipage}[b]{\linewidth}\centering
\textbf{\hl{Main reconstruction challenges}}
\end{minipage} & \begin{minipage}[b]{\linewidth}\centering
\textbf{\hl{Relevance to comparative evaluation}}
\end{minipage} \\
\midrule\noalign{}
\endfirsthead
\toprule\noalign{}
\begin{minipage}[b]{\linewidth}\centering
\textbf{\hl{Case ID}}
\end{minipage} & \begin{minipage}[b]{\linewidth}\centering
\textbf{\hl{Scene}}
\end{minipage} & \begin{minipage}[b]{\linewidth}\centering
\textbf{\hl{Architectural characteristics}}
\end{minipage} & \begin{minipage}[b]{\linewidth}\centering
\textbf{\hl{Main reconstruction challenges}}
\end{minipage} & \begin{minipage}[b]{\linewidth}\centering
\textbf{\hl{Relevance to comparative evaluation}}
\end{minipage} \\
\midrule\noalign{}
\endhead
\bottomrule\noalign{}
\endlastfoot
\hl{DDT-01} & \hl{Chaoli Drum Tower} & \hl{Multi-layered timber tower
with complex roof hierarchy} & \hl{Large viewpoint discontinuities;
unstable test-view construction; shadowed eave regions} & \hl{Testing
robustness under difficult acquisition conditions} \\
\hl{DDT-02} & \hl{Congjiang Drum Tower} & \hl{Distinct external geometry
with relatively regular overall form} & \hl{Repetitive timber textures;
local detail preservation} & \hl{Comparing rendering fidelity and
texture reconstruction} \\
\hl{DDT-03} & \hl{Zeli Drum Tower} & \hl{Strongly layered roof structure
with repeated eave patterns} & \hl{Repeated roof geometry; depth
ambiguity between eave layers} & \hl{Evaluating structural consistency
in complex roof regions} \\
\hl{DDT-04} & \hl{Zengchong Drum Tower} & \hl{Dense multi-eave geometry
with rich timber details} & \hl{Dense components; high-frequency detail
preservation; shadowed regions} & \hl{Testing detail reconstruction and
rendering robustness} \\
\end{longtable}

\subsection{\texorpdfstring{\hl{3.1 UAV Data Acquisition and Image
Corpus
Construction}}{3.1 UAV Data Acquisition and Image Corpus Construction}}\label{uav-data-acquisition-and-image-corpus-construction}

UAV images were acquired using a DJI Air 2S platform equipped with its
integrated camera sensor. The acquisition strategy was designed to
record the external morphology of the selected Dong drum towers as
completely as possible, including roof-crown structures, layered eaves,
timber façades, and immediate spatial surroundings. To balance global
coverage and local detail recording, the field campaign combined
sequential video recording for overall scene capture with still-image
acquisition for local observations. Multi-height and multi-angle views
were collected where site conditions allowed.

The acquisition parameters and image corpus information are summarized
in Table 3. In this study, overlap refers to the approximate average
overlap between adjacent UAV images during acquisition. Because the
source records provide only average overlap values rather than separate
forward and side overlap, the values are reported as nominal overlap
indicators. They are used to describe acquisition density and image
redundancy, not as independently measured forward- and side-overlap
parameters.

For Congjiang and Zengchong, twenty ground control points were measured
using Qianxun SR6 Plus differential GPS, with reported positioning error
below 1 cm. For Chaoli and Zeli, no GCP/RTK information is available;
therefore, scale reliability in these scenes should be interpreted with
caution and, where possible, checked against measurable architectural
dimensions such as total height, base width, column spacing, or eave
width. Weather and lighting conditions were not systematically recorded
during the field campaign; this limitation is acknowledged in the later
discussion.

All images were standardized to a resolution of 1920 × 1080 pixels
before neural reconstruction. The empirical corpus was constructed
through video-frame extraction, image screening, camera-pose estimation,
and scene-level data organization. The train/test split followed an
LLFF-style protocol, in which one image was selected as a test image at
regular intervals where applicable. The Chaoli scene has no independent
test image under the current split; therefore, its reported quantitative
results are interpreted as reference values rather than fully
independent test-set performance.

\textbf{Table 3. UAV acquisition parameters and empirical image corpus
of the four Dong drum tower scenes}

\begin{longtable}[]{@{}
  >{\centering\arraybackslash}p{(\linewidth - 22\tabcolsep) * \real{0.0593}}
  >{\centering\arraybackslash}p{(\linewidth - 22\tabcolsep) * \real{0.1041}}
  >{\centering\arraybackslash}p{(\linewidth - 22\tabcolsep) * \real{0.0805}}
  >{\centering\arraybackslash}p{(\linewidth - 22\tabcolsep) * \real{0.1004}}
  >{\raggedleft\arraybackslash}p{(\linewidth - 22\tabcolsep) * \real{0.0781}}
  >{\centering\arraybackslash}p{(\linewidth - 22\tabcolsep) * \real{0.1154}}
  >{\raggedleft\arraybackslash}p{(\linewidth - 22\tabcolsep) * \real{0.0585}}
  >{\raggedleft\arraybackslash}p{(\linewidth - 22\tabcolsep) * \real{0.0807}}
  >{\centering\arraybackslash}p{(\linewidth - 22\tabcolsep) * \real{0.0883}}
  >{\centering\arraybackslash}p{(\linewidth - 22\tabcolsep) * \real{0.1016}}
  >{\raggedleft\arraybackslash}p{(\linewidth - 22\tabcolsep) * \real{0.0805}}
  >{\raggedleft\arraybackslash}p{(\linewidth - 22\tabcolsep) * \real{0.0528}}@{}}
\toprule\noalign{}
\begin{minipage}[b]{\linewidth}\centering
\textbf{Case ID}
\end{minipage} & \begin{minipage}[b]{\linewidth}\centering
\textbf{Scene}
\end{minipage} & \begin{minipage}[b]{\linewidth}\centering
\textbf{Flight heights (m)}
\end{minipage} & \begin{minipage}[b]{\linewidth}\centering
\textbf{Image source}
\end{minipage} & \begin{minipage}[b]{\linewidth}\centering
\textbf{Overlap (\%)}
\end{minipage} & \begin{minipage}[b]{\linewidth}\centering
\textbf{GCP/RTK information}
\end{minipage} & \begin{minipage}[b]{\linewidth}\centering
\textbf{GSD (cm)}
\end{minipage} & \begin{minipage}[b]{\linewidth}\centering
\textbf{Final image corpus}
\end{minipage} & \begin{minipage}[b]{\linewidth}\centering
\textbf{Resolution}
\end{minipage} & \begin{minipage}[b]{\linewidth}\centering
\textbf{Acquisition mode}
\end{minipage} & \begin{minipage}[b]{\linewidth}\centering
\textbf{Training set}
\end{minipage} & \begin{minipage}[b]{\linewidth}\centering
\textbf{Test set}
\end{minipage} \\
\midrule\noalign{}
\endhead
\bottomrule\noalign{}
\endlastfoot
DDT-01 & Chaoli Drum Tower & 80 / 70 / 60 / 50 & 33 still images; 191
video frames & 75 & Not recorded & 4.35 & 490 & 1920 × 1080 & Circular
acquisition & 490 & 0* \\
DDT-02 & Congjiang Drum Tower & 70 / 60 / 40 & 96 still images; 94 video
frames & 70 & 20 GCPs; RTK error \textless{} 1 cm & 3.27 & 555 & 1920 ×
1080 & Circular acquisition & 494 & 61 \\
DDT-03 & Zeli Drum Tower & 50 / 40 / 30 & 99 still images; 115 video
frames & 80 & Not recorded & 2.96 & 588 & 1920 × 1080 & Circular
acquisition & 514 & 74 \\
DDT-04 & Zengchong Drum Tower & 75 / 50 / 40 / 30 & 86 still images; 159
video frames & 80 & 20 GCPs; RTK error \textless{} 1 cm & 3.54 & 519 &
1920 × 1080 & Circular acquisition & 458 & 61 \\
\end{longtable}

\begin{quote}
\textbf{Note:} GSD denotes ground sampling distance. Still images refer
to directly captured UAV photographs, whereas video frames refer to
images extracted from UAV video sequences for photogrammetric and neural
reconstruction. Overlap values indicate approximate average image
overlap during acquisition. The final image corpus refers to the
standardized images retained after frame extraction, screening, and
scene-level organization. The Chaoli scene has no independent test image
under the current LLFF-style split.
\end{quote}

\section{3.2 Empirical Point-Cloud and Reconstruction Corpus
Construction}\label{empirical-point-cloud-and-reconstruction-corpus-construction}

The UAV data were further processed to construct an empirical corpus for
neural rendering and reconstruction experiments. The workflow included
frame extraction from UAV videos, image quality screening, camera pose
estimation, sparse reconstruction, dense point-cloud generation where
available, and image normalization for different reconstruction methods.
The corpus includes raw or extracted UAV images, image metadata,
SfM-based sparse point clouds, dense point clouds, reconstruction
outputs from four neural methods, and selected rendered views for
quantitative and qualitative evaluation.

The UAV-derived point-cloud statistics are summarized in Table 4. These
point clouds were generated during the photogrammetric reconstruction
workflow and provide scene-level information on data completeness,
missing regions, and reconstruction difficulty. They are used as spatial
references for corpus description, but not as independent geometric
ground truth for metric accuracy evaluation.

\textbf{\hl{Table 4. Summary statistics of the UAV-derived point-cloud
datasets}}

\begin{longtable}[]{@{}
  >{\centering\arraybackslash}p{(\linewidth - 16\tabcolsep) * \real{0.0689}}
  >{\centering\arraybackslash}p{(\linewidth - 16\tabcolsep) * \real{0.1387}}
  >{\centering\arraybackslash}p{(\linewidth - 16\tabcolsep) * \real{0.0905}}
  >{\centering\arraybackslash}p{(\linewidth - 16\tabcolsep) * \real{0.0947}}
  >{\centering\arraybackslash}p{(\linewidth - 16\tabcolsep) * \real{0.1023}}
  >{\centering\arraybackslash}p{(\linewidth - 16\tabcolsep) * \real{0.1051}}
  >{\centering\arraybackslash}p{(\linewidth - 16\tabcolsep) * \real{0.0774}}
  >{\centering\arraybackslash}p{(\linewidth - 16\tabcolsep) * \real{0.1401}}
  >{\centering\arraybackslash}p{(\linewidth - 16\tabcolsep) * \real{0.1822}}@{}}
\toprule\noalign{}
\begin{minipage}[b]{\linewidth}\centering
\textbf{Case ID}
\end{minipage} & \begin{minipage}[b]{\linewidth}\centering
\textbf{Scene}
\end{minipage} & \begin{minipage}[b]{\linewidth}\centering
\textbf{Valid images}
\end{minipage} & \begin{minipage}[b]{\linewidth}\centering
\textbf{Sparse points}
\end{minipage} & \begin{minipage}[b]{\linewidth}\centering
\textbf{Dense points}
\end{minipage} & \begin{minipage}[b]{\linewidth}\centering
\textbf{Point density}
\end{minipage} & \begin{minipage}[b]{\linewidth}\centering
\textbf{Noise level}
\end{minipage} & \begin{minipage}[b]{\linewidth}\centering
\textbf{Estimated coverage (\%)}
\end{minipage} & \begin{minipage}[b]{\linewidth}\centering
\textbf{Major missing regions}
\end{minipage} \\
\midrule\noalign{}
\endhead
\bottomrule\noalign{}
\endlastfoot
DDT-01 & Chaoli Drum Tower & 490 & 287,190 & 14,631,743 & Medium & Low &
85 & Inner base and lower rear side \\
DDT-02 & Congjiang Drum Tower & 555 & 50,648 & 13,801,780 & Relatively
high & Low & 90 & Tower tip and internal structure \\
DDT-03 & Zeli Drum Tower & 588 & 364,582 & 39,357,222 & Relatively high
& Low & 88 & Shadowed areas beneath the eaves \\
DDT-04 & Zengchong Drum Tower & 519 & 72,546 & 19,631,785 & Medium & Low
& 87 & Inner side of dense multi-layered eaves \\
\end{longtable}

\begin{quote}
\textbf{Note:} Valid images refer to the images used for point-cloud
reconstruction after data screening. Sparse and dense points were
generated during the photogrammetric reconstruction workflow. Point
density and noise level are reported as relative quality indicators
after point-cloud filtering. Estimated coverage indicates the
approximate completeness of the reconstructed external structure.
\end{quote}
